% 中文摘要

工业机器人编码器校准偏差会引入系统性误差，同时腐蚀控制回路（通过错误的雅可比矩阵计算）和状态观测（通过错误的正运动学），使机器人控制问题从完全可观测的马尔可夫决策过程（MDP）退化为部分可观测马尔可夫决策过程（POMDP）。本文研究在未知编码器偏差下，如何在不停机重新标定的前提下保持鲁棒的操作性能。

首先，我们将编码器偏差问题形式化为POMDP，其中偏差向量作为episode级隐变量使真实关节状态不可观测。通过可观测性分析，我们证明标准本体感觉观测不足以区分不同的偏差条件，而引入外部末端执行器位置信号可以恢复近似可观测性——有效地将POMDP还原为可用标准强化学习算法求解的MDP。

基于上述分析，我们提出一个两层系统：（1）\textbf{鲁棒任务策略}，采用RLPD在随机编码器偏差下训练，在模拟抓取任务上全偏差范围$[0, 0.3]$\,rad内达到92--100\%成功率；（2）\textbf{安全备份策略}，通过主动避碰保护人在回路训练过程中的人类操作者安全。域随机化技术弥合了备份策略的仿真-真机差距。

消融实验证实外部位置信号对偏差鲁棒性能是必要的，验证了理论可观测性分析。我们还提出了面向Franka Panda真机的部署方案，包括双注入点故障注入机制，可在物理硬件上忠实复现编码器偏差因果链。

\vspace{1em}
\noindent\textbf{关键词：}部分可观测马尔可夫决策过程；强化学习；容错控制；编码器偏差；机器人操作
